% ============================================================
%  PART V --- Settings
% ============================================================
\part{Settings}

% ── Chapter 20: Settings Overview ───────────────────────────
\chapter{Settings --- Overview}
\label{ch:settings}

Open Settings by tapping \iconlabel{settings-iec}{Settings} in the Bottom Bar.
Desktop and iOS configuration is arranged across ten tabs. Android adds an
eleventh \textbf{Android} tab after System.
Changes take effect immediately unless a restart is explicitly noted.

\screenshot{settings-overview}{Settings panel with the platform-appropriate tabs visible in the tab bar}{}

\rowcolors{2}{LtGray}{white}
\begin{longtable}{C{0.6cm}L{3.0cm}L{6.8cm}L{3.3cm}}
\toprule
\rowcolor{Navy}
\color{white}\sffamily\bfseries \# &
\color{white}\sffamily\bfseries Tab name &
\color{white}\sffamily\bfseries Primary function &
\color{white}\sffamily\bfseries Restart needed? \\
\midrule
0 & Main           & Window mode, language, UI scale, comfort preset, launcher grid, coordinate format & Language and virtual keyboard only \\
1 & Top Bar        & Drag-and-drop layout of the Top Bar instrument strip & No \\
2 & Bottom Bar     & Drag-and-drop layout of the Bottom Bar action strip & No \\
3 & Widgets        & Define and edit reusable data widget definitions & No \\
4 & Comfort        & Select and customise visual comfort presets & No \\
5 & Connection     & Signal~K server URL, authentication token, mDNS discovery & No \\
6 & Signal~K       & Map semantic names to actual Signal~K paths & No \\
7 & Units          & Choose display units per measurement category & No \\
8 & Applications   & Manage the launcher: pages, tiles, folders, app details & No \\
9 & System         & Diagnostics, logging, performance, configuration management & No (some actions restart) \\
10 & Android       & Optional Home role, native launcher application selection, soft navigation controls & No (Android only) \\
\bottomrule
\end{longtable}

On Android, the Android tab lists only launchable \code{MAIN + LAUNCHER}
activities. Selecting an app adds it to the shared Applications palette without
requesting broad package visibility. Native app identity, icon, package, and
activity are Android-managed. Deselecting an app removes its launcher-page
references. FairWindSK never becomes the default Home app automatically, and the
ordinary launcher intent remains available. Soft Back, Home, and Recents controls
appear only when Android reports that the corresponding hardware keys are absent;
Recents contains only native activities launched through FairWindSK.

% ── Chapter 21: Tab 0 Main ──────────────────────────────────
\chapter{Settings --- Main}
\label{ch:settings-main}
\settingstab{0}{Main}

The Main tab controls fundamental runtime behaviour: window display mode, language, UI element
scale, active comfort preset, launcher grid dimensions, and coordinate format.

\screenshot{settings-main}{Settings --- Main tab: window mode, language, UI scale, comfort preset, launcher grid, coordinate format}{}

\section{Window Mode}

\rowcolors{2}{LtGray}{white}
\begin{longtable}{L{2.4cm}L{5.6cm}L{5.7cm}}
\toprule
\rowcolor{Navy}
\color{white}\sffamily\bfseries Mode &
\color{white}\sffamily\bfseries Behaviour &
\color{white}\sffamily\bfseries Notes \\
\midrule
Windowed    & Resizable desktop window at the configured Left/Top/Width/Height. &
  Development, testing, multi-monitor setups. \\
Centred     & Centred window of the configured Width and Height. Position fields disabled. &
  Preferred windowed mode without an exact position. \\
Maximised   & Full desktop work area. &
  On Raspberry Pi~OS, uses the screen work area geometry directly,
  bypassing the window manager. \\
Full Screen & Frameless, always-on-top kiosk window. &
  Recommended for dedicated helm displays.
  On labwc desktops, set \code{autohide=true}
  in \code{wf-panel-pi.ini} if the panel overlays FairWindSK. \\
\bottomrule
\end{longtable}

\section{Language}

FairWindSK ships with two fully supported interface languages.
The language setting affects all native shell text --- bar labels, settings pages,
drawers, warnings, and status messages --- as well as the \code{Accept-Language}
header sent to embedded web content.

\begin{itemize}
  \item \textbf{System default} --- follows the OS locale when FairWindSK supports
        that language; falls back to English for every unsupported OS language.
  \item \textbf{English} --- forces English regardless of OS locale.
  \item \textbf{Italiano} --- forces Italian. All native shell text is displayed in Italian.
\end{itemize}

\subsection{How localisation works}

The \code{Localization} module normalises the selected language, chooses the
effective \code{QLocale}, and builds the web \code{Accept-Language} header.
The Qt translator (\code{.qm} resource) is loaded during startup, so native widgets
and embedded web views start with the same language and culture.
Translation source files live in \code{resources/i18n/} (e.g.\
\code{fairwindsk\_it.ts} for Italian).

\subsection{Adding or improving a language}

See \code{docs/multilingual.md} in the repository for the full contributor workflow:
adding translatable strings with \code{tr(\ldots)}, updating \code{.ts} files with
Qt Linguist, registering a new language in \code{Localization.cpp} and
\code{ui/settings/Main.cpp}, and updating \code{CMakeLists.txt}.

\begin{warningbox}
A restart is required after changing the language.
FairWindSK shows an informational drawer when the language is changed.
Without a restart, native Qt widgets and the embedded web view may use different
languages.
\end{warningbox}

\section{UI Scale Preset}

Four presets: Small, Normal (default), Large, Extra Large.
When \key{Auto Scale} is ticked, FairWindSK selects the preset based on display resolution
and DPI.

\section{Comfort View Preset and Automatic Switching}

Selects the active comfort preset directly from the Main tab.
When \key{Auto Comfort View} is enabled, FairWindSK switches presets automatically
using Signal~K sun-position data (\skpath{environment.sun}).
Both conditions must be met for auto mode to be available:

\begin{enumerate}
  \item The \skpath{environment.sun} path must be configured in the Signal~K tab.
  \item The Signal~K server must be connected and delivering values on that path.
\end{enumerate}

\section{Launcher Grid}

Two spin boxes set the number of rows and columns of app tiles per launcher page.
Fewer rows and columns produce larger tiles that are easier to tap.
More produce more apps visible per page.

\section{Coordinate Format}

Three formats are available throughout the interface (Top Bar position widget, POB bar,
Anchor bar, My~Data waypoints):

\rowcolors{2}{LtGray}{white}
\begin{longtable}{L{3.2cm}L{4.8cm}L{5.7cm}}
\toprule
\rowcolor{Navy}
\color{white}\sffamily\bfseries Format &
\color{white}\sffamily\bfseries Example &
\color{white}\sffamily\bfseries When to use \\
\midrule
DD.DDDDD\textdegree{} & 48.85341\textdegree{} N & GPS logs, digital charting software. \\
DD\textdegree{}MM.MMM' & 48\textdegree{}51.205' N & Standard nautical format. Recommended for most sailors. \\
DD\textdegree{}MM'SS.S'' & 48\textdegree{}51'12.3'' N & Traditional celestial navigation format. \\
\bottomrule
\end{longtable}

\section{Virtual Keyboard}

Enables the Qt virtual keyboard for touchscreen-only deployments.
A restart is required after toggling.
On Raspberry~Pi~OS, the startup helper script applies the correct environment variable
(\code{QT\_IM\_MODULE=qtvirtualkeyboard}) automatically.

% ── Chapter 22: Tabs 1 and 2 --- Bar Editors ──────────────────
\chapter{Settings --- Top Bar and Bottom Bar Editors}
\label{ch:settings-bars}
\settingstab{1}{Top Bar} \settingstab{2}{Bottom Bar}

Tabs~1 and~2 share the same visual drag-and-drop editor --- one for the Top~Bar and one
for the Bottom~Bar.

\screenshot{settings-topbar-editor}{Settings --- Top Bar editor: live preview strip, editing toolbar, and widget palette}{}

\section{Live Preview Strip}

The preview at the top renders a full-width simulation of the bar using the current comfort
preset colours and real widget data.
Tap any item in the preview to select it.

\section{Editing Toolbar Buttons}

These are the exact icons defined in \code{BarLayoutSettings.cpp}.

\rowcolors{2}{LtGray}{white}
\begin{longtable}{C{1.5cm}L{4.0cm}L{8.2cm}}
\toprule
\rowcolor{Navy}
\color{white}\sffamily\bfseries Icon &
\color{white}\sffamily\bfseries Button &
\color{white}\sffamily\bfseries Action \\
\midrule
\iconlg{layout-horizontal-minimize} & Minimize width  &
  Widget uses only the space its content requires. \\
\iconlg{layout-horizontal-maximize} & Maximize width  &
  Widget expands to fill all available space horizontally. \\
\iconlg{layout-vertical-minimize}   & Minimize height &
  Widget uses the standard bar height. \\
\iconlg{layout-vertical-maximize}   & Maximize height &
  Widget uses the full bar height (both header and value rows fully visible). \\
\iconlg{arrow-left-google}          & Move left       &
  Moves the selected widget one position to the left. \\
\iconlg{arrow-right-google}         & Move right      &
  Moves the selected widget one position to the right. \\
\iconlg{delete-google}              & Remove selected &
  Returns the selected widget to the palette. \\
\iconlg{refresh-google}             & Reset defaults  &
  Restores the bar to its built-in default layout. \\
\bottomrule
\end{longtable}

\section{Display Toggles}

With a widget selected, four checkboxes control what is shown:
\key{Show Icon}, \key{Show Text} (label), \key{Show Units}, \key{Show Trend} (arrow).

\section{Widget Palette}

The palette below the toolbar lists all available items.
\textbf{Data widgets} are backed by Signal~K paths.
\textbf{Fixed shell controls} --- Apps, POB, Autopilot, Anchor, Alarms, My~Data, Settings,
Signal~K status, Clock --- can be moved between the Top~Bar and Bottom~Bar.
\textbf{Separators} (\iconlg{layout-separator}) and \textbf{Stretches}
(\iconlg{layout-elastic-extender}) are invisible layout helpers.

% ── Chapter 23: Tab 3 --- Widgets ─────────────────────────────
\chapter{Settings --- Widgets}
\label{ch:settings-widgets}
\settingstab{3}{Widgets}

The Widgets tab manages the library of data widget definitions used by both bars.

\screenshot{settings-widgets}{Settings --- Widgets tab with widget list and definition fields}{}

\section{Widget Definition Fields}

\rowcolors{2}{LtGray}{white}
\begin{longtable}{L{2.8cm}L{10.9cm}}
\toprule
\rowcolor{Navy}
\color{white}\sffamily\bfseries Field &
\color{white}\sffamily\bfseries Purpose \\
\midrule
id            & Unique identifier used internally and in bar layout entries. \\
name          & Human-readable label shown in bars and this editor. \\
icon          & Optional SVG icon path. Accepts Qt resource paths or filesystem paths. \\
type          & Rendering format: \code{numerical}, \code{gauge}, \code{position},
                \code{datetime}, or \code{waypoint}. \\
signalKPath   & Full dot-separated Signal~K data path. \\
sourceUnit    & Unit in which Signal~K delivers the value (e.g.\ \code{m/s}, \code{rad}). \\
defaultUnit   & Unit to display in (e.g.\ \code{kn}, \code{deg}).
                Overridden by active unit preferences. \\
updatePolicy  & \code{instant} (default) or \code{fixed}. \\
period        & Target update period in milliseconds (default 1000). \\
minPeriod     & Minimum update period in milliseconds (default 200). \\
\bottomrule
\end{longtable}

% ── Chapter 24: Tab 4 --- Comfort ─────────────────────────────
\chapter{Settings --- Comfort}
\label{ch:settings-comfort}
\settingstab{4}{Comfort}

The Comfort tab is the visual theme editor for all six built-in comfort presets.

\screenshot{settings-comfort}{Settings --- Comfort tab: preset selector, colour swatches, and background image section}{}

\section{The Six Comfort Presets}

\rowcolors{2}{LtGray}{white}
\begin{longtable}{C{1.5cm}L{2.0cm}L{4.2cm}L{6.0cm}}
\toprule
\rowcolor{Navy}
\color{white}\sffamily\bfseries Icon &
\color{white}\sffamily\bfseries Preset &
\color{white}\sffamily\bfseries Conditions &
\color{white}\sffamily\bfseries Key characteristics \\
\midrule
\iconlg{comfort-default} & Default & General fallback &
  Balanced mid-contrast. Usable in most conditions. \\
\iconlg{comfort-dawn}    & Dawn    & Around sunrise &
  Warm tones, slightly reduced brightness. Good contrast. \\
\iconlg{comfort-day}     & Day     & Full daylight, direct sun &
  Maximum brightness and contrast. For open-ocean sailing. \\
\iconlg{comfort-sunset}  & Sunset  & Around sunset &
  Warmer, transitional brightness. \\
\iconlg{comfort-dusk}    & Dusk    & Twilight &
  Reduced brightness, begins night-vision preservation. \\
\iconlg{comfort-night}   & Night   & Full dark, offshore watches &
  Minimum brightness, dark backgrounds, red accents.
  Preserves night vision. Alarms remain high-contrast. \\
\bottomrule
\end{longtable}

\section{Action Buttons}

\rowcolors{2}{LtGray}{white}
\begin{longtable}{C{1.5cm}L{3.0cm}L{9.2cm}}
\toprule
\rowcolor{Navy}
\color{white}\sffamily\bfseries Icon &
\color{white}\sffamily\bfseries Button &
\color{white}\sffamily\bfseries Action \\
\midrule
\iconlg{anchor-reset}     & Reset       &
  Clears all custom colour overrides, background images, and QSS for the selected preset. \\
\iconlg{delete-google}    & Reset All   &
  Clears customisation for all six presets. \\
\iconlg{route-export-iec} & Export QSS  &
  Saves the selected preset's effective stylesheet to a \code{.qss} file. \\
\iconlg{route-import-iec} & Import QSS  &
  Loads a \code{.qss} file as the custom stylesheet for the selected preset. \\
\bottomrule
\end{longtable}

\section{Colour Override Properties}

Ten colour properties are customisable per preset:

\rowcolors{2}{LtGray}{white}
\begin{longtable}{L{3.6cm}L{10.1cm}}
\toprule
\rowcolor{Navy}
\color{white}\sffamily\bfseries Property &
\color{white}\sffamily\bfseries What it controls \\
\midrule
Window background  & Background colour of the Application Area and settings pages. \\
Text               & Primary text colour throughout the interface. \\
Button background  & Background of buttons and tool buttons in bars and drawers. \\
Button text        & Text and icon tint on buttons. \\
Border             & Borders, dividers, and frame outlines. \\
Accent             & Primary accent colour for highlights and selected states. \\
Accent text        & Text colour on accent-coloured backgrounds. \\
Icon               & Default tint for monochrome SVG icons. Set to red for Night preset to aid night vision. \\
Scroll track       & Background of scroll bar tracks. \\
Scroll knob        & Colour of the scroll bar handle. \\
\bottomrule
\end{longtable}

\begin{warningbox}
Never configure a preset that makes alarm text, POB information, or safety-critical
readouts hard to read.
Readability under all conditions takes precedence over aesthetics.
\end{warningbox}

% ── Chapter 25: Tab 5 --- Connection ──────────────────────────
\chapter{Settings --- Connection}
\label{ch:settings-connection}
\settingstab{5}{Connection}

\screenshot{settings-connection}{Settings --- Connection tab: server URL field, status line, action buttons, and console log}{}

\section{Server URL Field}

An editable combo box pre-populated with the configured URL, \code{http://localhost:3000},
and \code{https://demo.signalk.org}.
FairWindSK normalises the URL: prepends \code{http://} if no scheme is present,
strips trailing slashes, and validates that only \code{http://} or \code{https://} are accepted.
Typed URLs are committed to the saved configuration only when you tap \key{Connect}.

\section{Status Line}

A single line below the URL field shows the current token and connection state using
three dot-separated fields: \textbf{State}, \textbf{Permission}, and \textbf{Expires}.

\section{Action Buttons}

\subsection{Connect and Pause}

Toggles the Signal~K connection.
\key{Connect} starts the connection; \key{Pause} suspends it without clearing the URL.

\subsection{Request Token}

Initiates the Signal~K read/write authentication flow:

\begin{enumerate}
  \item FairWindSK posts a new access request to \skpath{/signalk/v1/access/requests}.
  \item On HTTP~202 + \code{PENDING}, FairWindSK opens the Signal~K admin interface in the
        Bottom Bar drawer, navigated to the Security~> Access Requests page.
  \item FairWindSK polls the \code{href} every second until resolved.
  \item On \code{COMPLETED} with \code{APPROVED}, the token is stored in
        \code{fairwindsk.ini} and injected as the \code{JAUTHENTICATION} cookie in the
        embedded browser's cookie store.
\end{enumerate}

\subsection{Cancel}

Aborts the in-flight token request and clears the pending href.

\subsection{Remove Token}

Removes the stored token from memory, \code{fairwindsk.ini}, and the browser cookie store.
Triggers a reconnect in public (unauthenticated) mode.

\section{Console Log}

A read-only scrollable text area below the buttons shows connection events:
mDNS discovery messages, connection lifecycle events, and token request outcomes.

\section{mDNS Automatic Discovery}

On desktop platforms, FairWindSK scans the local network for Signal~K servers using
mDNS/Zeroconf on four service types:
\code{\_signalk-http.\_tcp}, \code{\_signalk-https.\_tcp}, \code{\_http.\_tcp},
and \code{\_https.\_tcp}.
Discovered servers are added to the URL combo box automatically.
mDNS is disabled on Android and iOS builds.

% ── Chapter 26: Tab 6 --- Signal K Paths ──────────────────────
\chapter{Settings --- Signal K Paths}
\label{ch:settings-signalk}
\settingstab{6}{Signal K Paths}

Maps the semantic names that FairWindSK uses internally to the actual Signal~K data paths
on your server.
Edit any field to change where FairWindSK reads or writes that value.
Changes take effect after a short debounce delay with no restart required.

\screenshot{settings-signalk}{Settings --- Signal K Paths tab with label and editable path fields}{}

\section{Key Path Mappings}

\rowcolors{2}{LtGray}{white}
\begin{longtable}{L{3.0cm}L{5.6cm}L{5.1cm}}
\toprule
\rowcolor{Navy}
\color{white}\sffamily\bfseries Key &
\color{white}\sffamily\bfseries Default Signal~K path &
\color{white}\sffamily\bfseries Feature \\
\midrule
\code{pos}              & \skpath{navigation.position}                              & Top~Bar position, POB position \\
\code{sog}              & \skpath{navigation.speedOverGround}                       & Top~Bar SOG widget \\
\code{cog}              & \skpath{navigation.courseOverGroundTrue}                   & Top~Bar COG widget \\
\code{hdg}              & \skpath{navigation.headingTrue}                           & Top~Bar HDG widget \\
\code{dpt}              & \skpath{environment.depth.belowTransducer}                & Top~Bar depth widget \\
\code{notifications.pob}& \skpath{notifications.mob}                                & POB alarm \\
\code{anchor.position}  & \skpath{navigation.anchor.position}                      & Anchor Watch \\
\code{anchor.radius}    & \skpath{navigation.anchor.currentRadius}                 & Anchor Watch \\
\code{anchor.max}       & \skpath{navigation.anchor.maxRadius}                     & Anchor Watch \\
\code{anchor.rode}      & \skpath{winches.windlass.rode}                           & Anchor rode display \\
\code{anchor.actions.drop}   & \skpath{plugins.anchoralarm.dropAnchor}          & Anchor Drop button \\
\code{anchor.actions.raise}  & \skpath{plugins.anchoralarm.raiseAnchor}         & Anchor Raise button \\
\code{anchor.actions.radius} & \skpath{plugins.anchoralarm.setRadius}           & Anchor Set Radius button \\
\code{anchor.actions.up}     & \skpath{plugins.windlassctl.up}                  & Anchor Up button \\
\code{anchor.actions.down}   & \skpath{plugins.windlassctl.down}                & Anchor Down button \\
\code{environment.sun}  & \skpath{environment.sun}                                  & Auto comfort preset switching \\
\code{autopilot.state}  & \skpath{steering.autopilot.state}                        & Autopilot bar mode display \\
\code{autopilot.mode}   & \skpath{steering.autopilot.mode}                         & Autopilot engage commands \\
\code{rsa}              & \skpath{steering.rudderAngle}                            & Autopilot rudder indicator \\
\bottomrule
\end{longtable}

% ── Chapter 27: Tab 7 --- Units ────────────────────────────────
\chapter{Settings --- Units}
\label{ch:settings-units}
\settingstab{7}{Units}

\screenshot{settings-units}{Settings --- Units tab: server preset name and per-category override rows}{}

\section{How Unit Preferences Work}

FairWindSK reads unit preferences from the connected Signal~K server.
Each category row shows: the category label, an override combo box (your local preference),
and the server's current value (read-only).
If you choose the same unit as the server, no override is stored.
If different, it is saved in \code{unitPreferences.overrides.CATEGORY} in
\code{fairwindsk.json} and takes precedence.

% ── Chapter 28: Tab 8 --- Applications ────────────────────────
\chapter{Settings --- Applications}
\label{ch:settings-apps}
\settingstab{8}{Applications}

A two-pane editor: the left pane manages launcher pages and folders; the right pane is
the application palette.

\screenshot{settings-apps}{Settings --- Applications tab: page tree (left) and app palette (right)}{}

\section{Page Tree Actions}

\rowcolors{2}{LtGray}{white}
\begin{longtable}{C{1.5cm}L{3.8cm}L{8.4cm}}
\toprule
\rowcolor{Navy}
\color{white}\sffamily\bfseries Icon &
\color{white}\sffamily\bfseries Action &
\color{white}\sffamily\bfseries What it does \\
\midrule
\iconlg{add-new-page}                    & Add page       & Creates a new empty launcher page. \\
\iconlg{show-page-details}               & Page details   & Opens the page details drawer (name, icon). \\
\iconlg{add-all-apps-to-current-page}    & Assign all apps & Adds all palette apps to the current page. \\
\iconlg{remove-all-apps-from-current-page}& Clear page    & Removes all tiles from the current page. \\
\iconlg{remove-all-apps-from-all-pages}  & Clear all pages & Removes all tiles from every page. \\
\bottomrule
\end{longtable}

\section{Palette Actions}

\rowcolors{2}{LtGray}{white}
\begin{longtable}{C{1.5cm}L{4.0cm}L{8.2cm}}
\toprule
\rowcolor{Navy}
\color{white}\sffamily\bfseries Icon &
\color{white}\sffamily\bfseries Action &
\color{white}\sffamily\bfseries What it does \\
\midrule
\iconlg{define-new-application-in-palette}  & Create app    & Opens the App Details drawer to create a new custom application. \\
\iconlg{show-details-selected-app-in-palette} & App details & Opens the App Details drawer to edit the selected application. \\
\iconlg{remove-selected-app-from-palette}   & Remove app    & Permanently removes the selected application from the palette. \\
\iconlg{add-selected-app-to-current-page}   & Add to page   & Adds the selected palette app to the current launcher page. \\
\bottomrule
\end{longtable}

\section{App Details Fields}

\rowcolors{2}{LtGray}{white}
\begin{longtable}{L{2.8cm}L{10.9cm}}
\toprule
\rowcolor{Navy}
\color{white}\sffamily\bfseries Field &
\color{white}\sffamily\bfseries Description \\
\midrule
Display name    & Name shown on the launcher tile and in the Top Bar when the app is active. \\
URL             & Full URL (\code{http://} or \code{https://}). \\
Icon path       & Qt resource path, filesystem path, or bundled icon path. \\
Web zoom level  & Percentage (default 100\%). Increase for larger buttons; decrease for more content. \\
Active          & Whether the tile is visible on the launcher. \\
Order           & Secondary sort hint within a page. \\
\bottomrule
\end{longtable}

% ── Chapter 29: Tab 9 --- System ──────────────────────────────
\chapter{Settings --- System}
\label{ch:settings-system}
\settingstab{9}{System}

The System tab is divided into two panes: \textbf{Performance} (left) and
\textbf{Diagnostics and Logging} (right).

\screenshot{settings-system-perf}{Settings --- System tab Performance pane: memory bars, per-core CPU, network, Raspberry Pi metrics}{}

\section{Performance Pane}

\subsection{Memory Readouts}

\begin{itemize}
  \item \textbf{Process memory} --- resident RAM used by FairWindSK
        (from \code{/proc/self/statm} on Linux; Mach task API on macOS).
  \item \textbf{Total memory} --- installed RAM
        (from \code{/proc/meminfo} on Linux; \code{sysctl hw.memsize} on macOS).
  \item \textbf{Memory usage bar} --- process memory as a percentage of total RAM.
\end{itemize}

\subsection{CPU Load}

A row of horizontal progress bars, one per CPU core, updated every second.
On Linux, data is read from \code{/proc/stat};
on macOS from the Mach \code{host\_processor\_info} API.

\subsection{Network Addresses}

Lists IPv4 addresses of all active, non-loopback network interfaces, refreshed every 10 seconds.
Useful for confirming which IP address other devices should use to reach this installation.

\subsection{Raspberry Pi Diagnostics}

When the connected Signal~K server exposes Raspberry~Pi metrics
(via the \code{signalk-server-raspberrypi} plugin on path \skpath{environment.rpi}),
a Raspberry~Pi group box appears showing:

\rowcolors{2}{LtGray}{white}
\begin{longtable}{L{3.6cm}L{10.1cm}}
\toprule
\rowcolor{Navy}
\color{white}\sffamily\bfseries Metric &
\color{white}\sffamily\bfseries Description \\
\midrule
CPU temperature   & SoC CPU temperature in \textdegree{}C. Above 80\textdegree{}C indicates thermal throttling. \\
GPU temperature   & VideoCore GPU temperature in \textdegree{}C. \\
CPU utilisation   & Overall CPU load as a percentage. \\
Memory utilisation& RAM usage as a percentage. \\
SD utilisation    & Storage usage on the SD card as a percentage. \\
\bottomrule
\end{longtable}

\section{Diagnostics and Logging Pane}

\screenshot{settings-system-log}{Settings --- System tab Diagnostics pane: log level, persistent logs, interaction history, paths, View Logs and View Reports buttons}{}

\subsection{Log Level}

\rowcolors{2}{LtGray}{white}
\begin{longtable}{L{2.4cm}L{4.6cm}L{6.7cm}}
\toprule
\rowcolor{Navy}
\color{white}\sffamily\bfseries Level &
\color{white}\sffamily\bfseries Messages included &
\color{white}\sffamily\bfseries When to use \\
\midrule
Off (default) & No messages. & Normal operation. Minimises SD card writes on Raspberry~Pi. \\
Critical      & Fatal errors only. & Minimal logging with crash detection. \\
Warning       & Critical + non-fatal anomalies. & Initial installation testing. \\
Info          & Warning + lifecycle events. & Commissioning and setup. \\
Debug         & Info + detailed internal state. & Troubleshooting. Significant overhead. \\
Full          & All of the above + verbose tracing. & Developer use only. Very high overhead. \\
\bottomrule
\end{longtable}

\subsection{Persistent Logs}

When enabled, FairWindSK writes a complete message log for each run to the per-user
application data directory, with the process start timestamp in the filename.

\subsection{Interaction History}

When enabled, FairWindSK writes a lightweight journal of operator navigation and control
actions alongside each run log.
Useful for post-incident review.

\subsection{View Logs and View Reports Buttons}

Open the in-application log explorer in the Bottom Bar horizontal drawer.
Log files appear on the left; parsed log entries appear on the right.
Crash-report bundles (created automatically when FairWindSK did not exit cleanly) are
accessible via \key{View Reports}.

\section{Configuration Management Buttons}

\rowcolors{2}{LtGray}{white}
\begin{longtable}{L{3.2cm}L{10.5cm}}
\toprule
\rowcolor{Navy}
\color{white}\sffamily\bfseries Button &
\color{white}\sffamily\bfseries Action \\
\midrule
Reset           &
  Discards unsaved changes and reloads the active configuration.
  FairWindSK keeps running. \\
Restore Defaults &
  Replaces all settings with the bundled defaults.
  All customisation is lost. Confirmation required. \\
Import          &
  Imports a \code{fairwindsk.json} file from the filesystem.
  Validates, atomically replaces the active configuration, and reloads immediately. \\
Export          &
  Copies the current \code{fairwindsk.json} to a chosen location. \\
\bottomrule
\end{longtable}

\begin{warningbox}
\key{Restore Defaults} is irreversible without a backup.
Export your configuration first if you may want to restore it.
\end{warningbox}

\section{System Control Buttons}

\rowcolors{2}{LtGray}{white}
\begin{longtable}{L{3.2cm}L{10.5cm}}
\toprule
\rowcolor{Navy}
\color{white}\sffamily\bfseries Button &
\color{white}\sffamily\bfseries Action \\
\midrule
Restart FairWindSK &
  Saves settings, then closes FairWindSK with exit code~1.
  The Raspberry~Pi~OS startup helper relaunches on exit code~1. \\
Quit FairWindSK &
  Saves settings and closes cleanly (exit code~0).
  The startup helper does \textbf{not} relaunch after a clean exit. \\
\bottomrule
\end{longtable}

\begin{warningbox}
Quitting FairWindSK on a dedicated helm display leaves the display blank.
Ensure another crew member is maintaining a safe watch before quitting.
\end{warningbox}
